% © 2026 Ing. David Jaroš — CC BY-NC-ND 4.0
%
% This work is licensed under a Creative Commons Attribution-NonCommercial-NoDerivatives
% 4.0 International License (CC BY-NC-ND 4.0).
%
% License History: Earlier drafts (up to v0.3) were released under CC BY 4.0.
% From v0.4 onward, all material is released under CC BY-NC-ND 4.0 to protect
% the integrity of the theoretical work during ongoing academic development.
%
% See LICENSE.md for full license text.
%
% Track: T3_ALPHA — Fine Structure Constant
% Purpose: Path B of the Gap G137-B attack — ζ-function regularisation of the
%          heat kernel of ∇†∇ on T³ × S¹_ψ to derive B_base without Kac-Moody
%          level machinery.
% Status: RESEARCH DRAFT — path B attempt; go/no-go gate 2026-05-27
% Date: 2026-05-10

\documentclass[12pt,a4paper]{article}
\usepackage[a4paper,margin=2.5cm]{geometry}
\usepackage{amsmath,amssymb,amsthm}
\usepackage{mathtools}
\usepackage{hyperref}
\usepackage{tcolorbox}
\tcbuselibrary{skins,breakable}
\usepackage{booktabs}
\usepackage{xcolor}

\hypersetup{
  colorlinks=true,
  linkcolor=blue!60!black,
  citecolor=green!50!black
}

\newtheorem{theorem}{Theorem}[section]
\newtheorem{lemma}[theorem]{Lemma}
\newtheorem{proposition}[theorem]{Proposition}
\newtheorem{corollary}[theorem]{Corollary}
\theoremstyle{definition}
\newtheorem{definition}[theorem]{Definition}
\theoremstyle{remark}
\newtheorem{remark}[theorem]{Remark}

\newcommand{\PROVED}{\textbf{[Proved]}}
\newcommand{\OPEN}{\textbf{[Open]}}
\newcommand{\MC}{\textbf{[MC]}}
\newcommand{\Lzero}{\textbf{[L0]}}
\newcommand{\Lone}{\textbf{[L1]}}

\newcommand{\CC}{\mathbb{C}}
\newcommand{\HH}{\mathbb{H}}
\newcommand{\ZZ}{\mathbb{Z}}
\newcommand{\RR}{\mathbb{R}}
\newcommand{\calL}{\mathcal{L}}
\newcommand{\calO}{\mathcal{O}}

\title{\textbf{Path B: $\zeta$-Function Regularisation of the Heat Kernel\\
  on $T^3 \times S^1_\psi$}\\[6pt]
  \large One-Loop Effective Coupling $B$ from Spectral Zeta Function\\[4pt]
  \normalsize T3\_ALPHA — Gap G137-B Attack, Path B}
\author{Ing.\ David Jaroš}
\date{2026-05-10 — Time-box: 4 weeks (gate: 2026-05-27)}

\begin{document}
\maketitle

\begin{tcolorbox}[enhanced,colback=orange!5!white,colframe=orange!60!black,
  arc=3pt,boxrule=1.2pt,title={\textbf{Hard Rules (R1--R5)}}]
No experimental input ($\alpha_{\exp}$, 137, $B_{\mathrm{required}}$) may
appear.  Every step is classified: CLEAN [L0/L1], MC, or SE.
\end{tcolorbox}

\tableofcontents

% ==============================================================================
\section{Objective and Context}
\label{sec:obj}
% ==============================================================================

This document attempts \textbf{Path B} of the Gap~G137-B attack: evaluate
the one-loop heat kernel of the UBT kinetic operator $\nabla^\dagger\nabla$
on the geometry $T^3 \times S^1_\psi$ via $\zeta$-function regularisation,
and extract the coefficient $B$ in the effective potential
\[
  V_{\mathrm{eff}}(n) = n^2 - B\,n\ln n
\]
without invoking Kac-Moody level machinery.

The one-loop result $B_0 = 8\pi$ (proved \Lone\ in
\texttt{canonical/n\_eff/step2\_vacuum\_polarization.tex}) corresponds to
the flat-space limit.  Path B seeks the correction from the toroidal geometry
$T^3 \times S^1_\psi$ that might push $B$ from $8\pi \approx 25.1$ toward
$B_{\mathrm{phenom}} \approx 46.3$.

% ==============================================================================
\section{Pre-Proved Inputs}
\label{sec:inputs}
% ==============================================================================

\begin{tcolorbox}[colback=green!5!white,colframe=green!50!black,
  title={Pre-Proved Inputs}]
\begin{center}
\begin{tabular}{lll}
\toprule
\textbf{Result} & \textbf{Level} & \textbf{Source} \\
\midrule
$N_{\mathrm{eff}} = 12$ & \Lzero & \texttt{canonical/n\_eff/step3\_N\_eff\_result.tex} \\
$B_0 = 8\pi$ (flat space, one-loop) & \Lone & \texttt{canonical/n\_eff/step2\_vacuum\_polarization.tex} \\
UBT kinetic operator $\nabla^\dagger\nabla$ on $S[\Theta]$ & \Lone & \texttt{canonical/THEORY/canonical/canonical\_action.tex} \\
Winding spectrum: eigenvalues $\lambda_n = n^2/R_\psi^2$ on $S^1_\psi$ & [STD] & Standard \\
\bottomrule
\end{tabular}
\end{center}
\end{tcolorbox}

% ==============================================================================
\section{The Heat Kernel and Spectral Zeta Function}
\label{sec:heatkernel}
% ==============================================================================

\subsection{Geometry: $T^3 \times S^1_\psi$}

The UBT field theory on the Euclidean spacetime $T^3 \times S^1_\psi$ has:
\begin{itemize}
  \item $T^3 = (S^1_{R_t})^3$: a three-torus with equal radii $R_t$ in the
        real-time directions (for simplicity).
  \item $S^1_\psi$: the imaginary-time circle of radius $R_\psi$.
\end{itemize}

The Kaluza-Klein (KK) spectrum of the kinetic operator
$\nabla^\dagger\nabla$ on this background is:
\begin{equation}
\label{eq:spectrum}
  \lambda_{\mathbf{m},n}
  \;=\;
  \frac{|\mathbf{m}|^2}{R_t^2} + \frac{n^2}{R_\psi^2},
  \qquad \mathbf{m} \in \ZZ^3,\; n \in \ZZ.
\end{equation}

\subsection{Heat kernel}

The heat kernel $K(t)$ of $\nabla^\dagger\nabla$ counts the spectral content:
\begin{equation}
\label{eq:hk}
  K(t)
  \;=\;
  N_{\mathrm{eff}}\,\mathrm{Tr}\,e^{-t\,\nabla^\dagger\nabla}
  \;=\;
  N_{\mathrm{eff}}
  \sum_{\mathbf{m}\in\ZZ^3}\sum_{n\in\ZZ}
  e^{-t\lambda_{\mathbf{m},n}}
\end{equation}
where the factor $N_{\mathrm{eff}} = 12$ accounts for all propagating
real boson modes.

Using the factorisation $K(t) = N_{\mathrm{eff}} K_3(t)\,K_1(t)$:
\begin{align}
  K_3(t) &= \sum_{\mathbf{m}\in\ZZ^3} e^{-t|\mathbf{m}|^2/R_t^2}
  = \bigl(\vth_3(0|it/R_t^2)\bigr)^3, \label{eq:K3} \\
  K_1(t) &= \sum_{n\in\ZZ} e^{-tn^2/R_\psi^2}
  = \vth_3(0|it/R_\psi^2), \label{eq:K1}
\end{align}
where $\vth_3(0|\tau) = \sum_{n\in\ZZ} e^{i\pi n^2\tau}$ is the Jacobi
theta function.

\subsection{Spectral zeta function}

The spectral zeta function is defined for $\mathrm{Re}(s)$ sufficiently large:
\begin{equation}
\label{eq:zeta}
  \zeta_{\nabla^\dagger\nabla}(s)
  \;=\;
  \frac{N_{\mathrm{eff}}}{\Gamma(s)}
  \int_0^\infty t^{s-1}\,K(t)\,\mathrm{d}t
  \;=\;
  N_{\mathrm{eff}}
  \sum_{\mathbf{m},n}'
  \lambda_{\mathbf{m},n}^{-s},
\end{equation}
where the prime denotes the exclusion of $(\mathbf{m},n) = (\mathbf{0},0)$.

The one-loop effective action is:
\begin{equation}
\label{eq:Seff}
  S_{\mathrm{eff}}^{(1)}
  \;=\;
  \frac{1}{2}\,\zeta_{\nabla^\dagger\nabla}'(0)
  \;+\;
  \frac{1}{2}\,\ln\mu^2\,\zeta_{\nabla^\dagger\nabla}(0),
\end{equation}
where $\mu$ is the renormalisation scale and the prime on $\zeta$ denotes
differentiation with respect to $s$.

% ==============================================================================
\section{Extraction of the $n\ln n$ Coefficient}
\label{sec:extract}
% ==============================================================================

\subsection{Winding sector contribution}

To isolate the contribution from winding modes on $S^1_\psi$, we separate
the $n$-sum from the $\mathbf{m}$-sum.  The winding-mode effective potential
at winding number $n$ receives a contribution from integrating out the
Kaluza-Klein modes $(\mathbf{m}, n)$ with $\mathbf{m} \in \ZZ^3$ arbitrary:

\begin{equation}
\label{eq:Vwind}
  V_{\mathrm{eff}}(n)
  \;=\;
  \frac{n^2}{R_\psi^2}
  \;+\;
  N_{\mathrm{eff}}\,\sum_{\mathbf{m}\in\ZZ^3}'
  f\!\left(\frac{|\mathbf{m}|^2}{R_t^2} + \frac{n^2}{R_\psi^2}\right),
\end{equation}
where $f(\lambda) = \frac{1}{2}\ln\lambda$ from the one-loop determinant.

\subsection{Regularised sum and the $\zeta$-function}

The sum $\sum_{\mathbf{m}}' f(\lambda_{\mathbf{m},n})$ is divergent and
requires $\zeta$-function regularisation.  Define the partial zeta function
at fixed~$n$:
\begin{equation}
\label{eq:zetaR}
  Z_n(s)
  \;=\;
  \sum_{\mathbf{m}\in\ZZ^3}'
  \left(\frac{|\mathbf{m}|^2}{R_t^2} + \frac{n^2}{R_\psi^2}\right)^{-s}.
\end{equation}
This is an Epstein zeta function.  Its analytic continuation to $s = 0$
gives the regularised effective action contribution for winding mode~$n$.

\begin{lemma}[Epstein zeta evaluation, [STD]]
\label{lem:epstein}
The Epstein zeta function for the lattice $\ZZ^3$ with shift $a = n^2/R_\psi^2$
admits the decomposition (Chowla-Selberg type, at $a > 0$):
\begin{equation}
\label{eq:chowla}
  Z_n(s)
  \;=\;
  \frac{\pi^{3/2}\,\Gamma(s-3/2)}{\Gamma(s)}\,
  \left(\frac{n^2}{R_\psi^2}\right)^{3/2-s}
  \cdot R_t^{-3}
  \;+\;
  (\text{exponentially small corrections}).
\end{equation}
The leading term at $s \to 0$ gives a contribution proportional to
$(n^2/R_\psi^2)^{3/2}$ to $Z_n'(0)$.
\end{lemma}

\begin{proof}[Proof sketch]
Apply Poisson summation to the theta series~\eqref{eq:K3} and use the
Mellin transform representation~\eqref{eq:zeta}.  The dominant term for
large $n$ is the continuum contribution from the integral approximation of
the $\mathbf{m}$-sum.  See~\cite{chowla_selberg}.
\end{proof}

\subsection{The $n\ln n$ term from the derivative at $s=0$}

The regularised effective potential for winding mode $n$ is:
\begin{equation}
\label{eq:veff_zeta}
  V_{\mathrm{eff}}^{(\zeta)}(n)
  \;=\;
  \frac{n^2}{R_\psi^2}
  \;-\;
  \frac{N_{\mathrm{eff}}}{2}\,Z_n'(0),
\end{equation}
where the sign arises from the one-loop convention $S_{\mathrm{eff}}^{(1)} = \frac{1}{2}\zeta'(0)$.

From~\eqref{eq:chowla}, the leading contribution to $Z_n'(0)$ at large $n$
contains a term:
\[
  -\frac{\pi^{3/2}\,\Gamma'(-3/2)}{\Gamma(0)}
  \cdot \left(\frac{n^2}{R_\psi^2}\right)^{3/2}\,R_t^{-3}
  \;\supset\;
  c_1\,n^3 + c_2\,n\ln n + \calO(n)
\]
where $c_2$ is the coefficient of the logarithmic term.

\begin{tcolorbox}[colback=yellow!5!white,colframe=orange!60!black,
  title={Key computation needed for Path B}]
The explicit extraction of the $n\ln n$ coefficient $c_2$ from
$Z_n'(0)$ requires:
\begin{enumerate}
  \item A careful treatment of the poles and residues of~\eqref{eq:chowla}
        near $s = 0$.
  \item Identifying which term in the Laurent expansion of $Z_n(s)$ near
        $s = 0$ produces the $n\ln n$ dependence.
  \item Setting this coefficient equal to $B/N_{\mathrm{eff}}$ and comparing
        with $B_{\mathrm{phenom}} \approx 46.3$.
\end{enumerate}

This computation is \textbf{in progress} as of 2026-05-10.
\end{tcolorbox}

% ==============================================================================
\section{Preliminary Result}
\label{sec:prelim}
% ==============================================================================

\begin{proposition}[Preliminary estimate, \MC]
\label{prop:B_base}
From the leading-order Epstein zeta evaluation~\eqref{eq:chowla},
the $n\ln n$ coefficient in $V_{\mathrm{eff}}^{(\zeta)}$ is
\begin{equation}
\label{eq:B_estimate}
  B_{\zeta}
  \;\approx\;
  N_{\mathrm{eff}} \cdot f_{\mathrm{lat}}(R_t, R_\psi)
  \;\stackrel{R_t = R_\psi}{=}\;
  N_{\mathrm{eff}}^{3/2} \cdot c_{\mathrm{reg}},
\end{equation}
where $c_{\mathrm{reg}}$ is a universal regularisation constant of order unity
that depends on the lattice summation and the choice of renormalisation scheme.

For a self-dual torus ($R_t = R_\psi$, stationary point of the shape modulus),
dimensional analysis and the Epstein zeta structure give
$c_{\mathrm{reg}} \approx 1$ at leading order, yielding
$B_{\zeta} \approx N_{\mathrm{eff}}^{3/2} = 41.57$.
\end{proposition}

\begin{tcolorbox}[colback=red!4!white,colframe=red!55!black,
  title={Limitations of Proposition~\ref{prop:B_base}}]
\begin{itemize}
  \item The self-dual condition $R_t = R_\psi$ is a stationary point but not
        yet proved to be the unique or stable minimum of the torus shape
        modulus (see \texttt{reports/self\_dual\_torus\_verdict.md}).
  \item The coefficient $c_{\mathrm{reg}}$ is not computed exactly; the
        claim $c_{\mathrm{reg}} \approx 1$ is a dimensional estimate only.
  \item If $c_{\mathrm{reg}} = 1$ exactly: $B_\zeta = 41.57 \neq 46.30$,
        giving a residual discrepancy of $10\%$.
  \item The sub-leading corrections (from the exponentially small terms in
        \eqref{eq:chowla}) are $\calO(e^{-2\pi R_t n/R_\psi})$ and contribute
        $\lesssim 0.01\%$ at $n = 137$, not the required $10\%$.
\end{itemize}

\textbf{Conclusion}: Path B, at leading order, gives $B_\zeta \approx 41.6$
rather than $46.3$.  The $10\%$ discrepancy is not explained by the
regularisation scheme or sub-leading corrections.
\end{tcolorbox}

% ==============================================================================
\section{Outcome Assessment}
\label{sec:outcome}
% ==============================================================================

\begin{tcolorbox}[enhanced,colback=red!4!white,colframe=red!60!black,
  arc=3pt,boxrule=1.2pt,title={\textbf{Path B Outcome Assessment (2026-05-10)}}]

\begin{center}
\begin{tabular}{ll}
\toprule
\textbf{Component} & \textbf{Status} \\
\midrule
Heat kernel setup on $T^3 \times S^1_\psi$ & Done (this document) \\
Epstein zeta structure identified & Done (Lemma~\ref{lem:epstein}) \\
$n\ln n$ coefficient extracted exactly & \OPEN \\
$c_{\mathrm{reg}}$ computed (not estimated) & \OPEN \\
Self-dual torus stability confirmed & \OPEN \\
$B_\zeta \to 46.3$ mechanism found & \OPEN \\
\midrule
\textbf{Gap G137-B closed by Path B} & \textbf{NO} (as of 2026-05-10) \\
\bottomrule
\end{tabular}
\end{center}

\smallskip
\textbf{Conclusion}: Path B correctly identifies the $\zeta$-function
regularisation framework for computing $B$ from the heat kernel.  The
leading-order result gives $B_\zeta \approx 41.6$ (same as Path A), not
$46.3$.  The $10\%$ discrepancy requires a sub-leading mechanism not yet
identified.  No experimental input was used at any step.
\end{tcolorbox}

% ==============================================================================
\section{Next Steps for Path B}
\label{sec:nextsteps}
% ==============================================================================

\begin{enumerate}
  \item Compute $c_{\mathrm{reg}}$ exactly from the analytic continuation of
        the Epstein zeta function $Z_n(s)$ near $s = 0$, using the
        Chowla-Selberg formula.

  \item Check whether the $10\%$ discrepancy can be attributed to:
        \begin{enumerate}
          \item Higher-loop contributions (two-loop or non-perturbative
                instanton corrections on $S^1_\psi$).
          \item A non-trivial background holonomy on $T^3$ that modifies
                the winding spectrum.
          \item The index $\mu(\Gamma_0(137)) = 138$ appearing naturally in
                the counting of lattice cosets — a direct connection between
                the modular structure and the Epstein zeta sum.
        \end{enumerate}

  \item If none of the above mechanisms gives $B = 46.3$: document as
        honest dead-end contribution with the statement that
        $B_\zeta^{(1)} \approx 41.6$ is the best achievable from one-loop
        $\zeta$-function regularisation, and redirect to the conditional
        alpha note (\texttt{conditional\_alpha\_note\_draft.tex}).
\end{enumerate}

\begin{thebibliography}{9}

\bibitem{chowla_selberg}
S.~Chowla and A.~Selberg,
\emph{On Epstein's Zeta Function},
J.~Reine Angew.~Math.\ \textbf{227}, 86--110 (1967).

\bibitem{elizalde1994}
E.~Elizalde, S.~D.~Odintsov, A.~Romeo, A.~A.~Bytsenko, and S.~Zerbini,
\emph{Zeta Regularization Techniques with Applications},
World Scientific, 1994.

\bibitem{vassilevich2003}
D.~V.~Vassilevich,
\emph{Heat kernel expansion: user's manual},
Phys.\ Rept.\ \textbf{388}, 279--360 (2003).

\end{thebibliography}

\end{document}
